\begin{exbox}[同一连续谱中的直接路径与共振路径：Fano 线形]
投影自能确定共振位置和宽度，却还未确定观测线形。
若探测器同时收集直接连续通道与离散共振衰变通道，就必须求两条路径的共同末态振幅。
取能量归一化连续态 $\langle E|E'\rangle=\delta(E-E')$，
\[
H=E_d|d\rangle\langle d|+\int E|E\rangle\langle E|\dd E
+\int[w(E)|E\rangle\langle d|+w(E)^*|d\rangle\langle E|]\dd E.
\]
$w$ 的单位为能量的平方根；态密度已包含在这个归一化中。
令激发算符作用于初态给
$D|i\rangle=\mu_d|d\rangle+\int M(E)|E\rangle\dd E$。
离散预解式为
\begin{align}
G_d^R(E)&=[E-E_d-\Sigma^R(E)]^{-1},
\label{eq:feqo-fano-selfenergy}\\
\Sigma^R(E)&=\mathop{\rm PV}\int\frac{|w(E')|^2}{E-E'}\dd E'
-\ii\pi|w(E)|^2\equiv\Delta_d(E)-\ii\Gamma_d(E)/2,
\label{eq:feqo-fano-width-shift}
\end{align}
其中 $\Gamma_d=2\pi|w|^2$ 为\emph{能量宽度}，衰减率是 $\Gamma_d/\hbar$。

把同一个预解式作用于激发源，先消去连续分量，再投影到出射连续态，得到
\[
T(E)=M(E)+w(E)G_d^R(E)
\left[\mu_d+\int\frac{w(E')^*M(E')}{E-E'+\ii0}\dd E'\right].
\]
这里的积分不能省掉：直接激发连续谱以后，它仍可散射到离散态再返回。
将其分成主值与在壳部分，令
\[
\mu_{\rm eff}(E)=\mu_d+\mathop{\rm PV}\int
\frac{w(E')^*M(E')}{E-E'}\dd E'.
\]
于是括号是 $\mu_{\rm eff}-\ii\pi w^*M$。
在单一连续通道、具有相应时间反演约束且跃迁矩阵元可取实数的模型中，
假定各慢变系数在共振宽度内可取常数，定义
\begin{equation}
\varepsilon=\frac{2(E-E_d-\Delta_d)}{\Gamma_d},
\qquad q_F=\frac{2w\mu_{\rm eff}}{M\Gamma_d}.
\label{eq:feqo-fano-q}
\end{equation}
将两条路径置于同一个分母，得到
\begin{align*}
T(E)&=M+\frac{w\mu_{\rm eff}-\ii\Gamma_dM/2}
{E-E_d-\Delta_d+\ii\Gamma_d/2}\\
&=M\frac{\varepsilon+q_F}{\varepsilon+\ii}.
\end{align*}
因此
\begin{equation}
|T(E)|^2=|M|^2\frac{(\varepsilon+q_F)^2}{1+\varepsilon^2}.
\label{eq:feqo-fano-lineshape}
\end{equation}
分子的虚项被直接通道的再散射精确抵消。
若只加 $M+\mu_dwG_d^R$，就会留下错误的分子常数 $+1$。
$|M|^2$ 是同一耦合连续通道的平滑背景，并非另一个互不干涉的连续谱。
有额外未相干收集的通道或真正损耗时，其强度应另外相加；
复杂相位通常使 $q_F$ 为复数，反共振未必是严格零点。
以上计算是 Fano 构型相互作用的单通道形式。\footnote{U.~Fano 的原始推导见
\url{https://doi.org/10.1103/PhysRev.124.1866}。}

对实数 $q_F\ne0$，求导得
\[
\frac{\dd}{\dd\varepsilon}\frac{(\varepsilon+q_F)^2}{1+\varepsilon^2}
=\frac{2(\varepsilon+q_F)(1-q_F\varepsilon)}{(1+\varepsilon^2)^2}.
\]
零点在 $\varepsilon=-q_F$，峰在 $\varepsilon=1/q_F$，峰高为背景的 $1+q_F^2$ 倍。
例如给定模型参数 $\Gamma_d=\SI{0.15}{\electronvolt}$、$q_F=2$，
零点在共振以下 $\SI{0.15}{\electronvolt}$，峰在共振以上
$\SI{0.0375}{\electronvolt}$，高度为背景五倍。
这组数只是线形模型的代入，不指定某一种材料都满足相同宽度。
\end{exbox}
